\documentclass[11pt]{article}
\usepackage[letterpaper,margin=0.82in]{geometry}
\usepackage{amsmath,amssymb,amsthm,mathtools}
\usepackage{graphicx}
\usepackage{microtype}
\usepackage{xcolor}
\usepackage[hidelinks]{hyperref}
\usepackage{enumitem}

\definecolor{deepblue}{RGB}{26,67,120}
\definecolor{softgray}{RGB}{245,247,250}
\newtheorem{theorem}{Theorem}
\newtheorem{lemma}[theorem]{Lemma}
\newtheorem{proposition}[theorem]{Proposition}
\theoremstyle{remark}
\newtheorem{remark}[theorem]{Remark}
\newcommand{\G}{\mathcal G}
\newcommand{\Hil}{\mathcal H}
\newcommand{\supp}{\operatorname{supp}}
\newcommand{\ran}{\operatorname{ran}}
\newcommand{\C}{\mathbb C}

\title{\vspace{-2.5em}\textbf{Fourier-algebra restriction to isotropy is always surjective}\\
\large A full answer to Remark 3.10 of arXiv:2411.14624}
\author{Candidate solution packet --- run \texttt{fa\_banach\_001}, lane 2}
\date{27 August 2026}

\begin{document}
\maketitle

\noindent
\colorbox{softgray}{\parbox{0.95\linewidth}{
\textbf{Status.} Candidate full solution, likely valid; human review is recommended.
For every locally compact Hausdorff groupoid with a left Haar system, restriction
from the groupoid Fourier algebra to the Fourier algebra of any isotropy group is
surjective. No countability, \'{e}tale, unimodularity, or right-invariance
hypothesis is required.}}

\section{Source question}

Joseph DeGaetani and Mahya Ghandehari, \emph{On the restriction maps of the
Fourier and Fourier-Stieltjes algebras over locally compact groupoids},
arXiv:2411.14624, prove the result for \'{e}tale groupoids in Theorem 3.7.
Remark 3.10 (printed page 13) asks whether that theorem is true in full
generality.  The exact source passage is reproduced below.

\begin{center}
\includegraphics[width=0.96\linewidth]{figures/open_problem_crop.png}
\end{center}

In symbols, if $\G$ is a locally compact Hausdorff groupoid with a left Haar
system $\{\lambda^v\}_{v\in\G^{(0)}}$, $u\in\G^{(0)}$, and
$H=\G_u^u$, the question is whether
\[
        R_u:A(\G)\longrightarrow A(H),\qquad R_u\varphi=\varphi|_H,
\]
is onto without assuming that $\G$ is \'{e}tale.

\section{Definitions and the source reduction}

Put $X=\G^u=r^{-1}(u)$ and $\mu=\lambda^u$.  The isotropy group
$H=\G_u^u$ acts on $X$ by left multiplication.  Fix a left Haar measure
$dh$ on $H$.  Write $\lambda$ for the left regular representation on
$L^2(H,dh)$,
\[
  (\lambda(k)\xi)(h)=\xi(k^{-1}h),
\]
and let $\sigma$ be the restriction to $H$ of the groupoid left-regular
representation,
\[
  (\sigma(k)g)(x)=g(k^{-1}x),\qquad g\in L^2(X,\mu).
\]
For a unitary representation $\pi$ of $H$, $A_\pi(H)$ denotes the closed
linear span in $B(H)$ of its coefficient functions. Thus
$A_\lambda(H)=A(H)$.

Two general steps already contained in the proof of Theorem 3.7 of the source
paper do not use etaleness:
\begin{equation}\label{eq:source-reduction}
       R_u(A(\G))=A_\sigma(H)\subseteq A(H).
\end{equation}
Indeed, each coefficient of $\sigma$ lifts to a coefficient of the groupoid
regular representation by the continuous-section extension theorem.  An
absolutely summable coefficient decomposition in $A_\sigma(H)$ then lifts
term by term, proving surjectivity onto $A_\sigma(H)$.  Conversely, restriction
of compactly supported groupoid coefficients lies in $C_c(H)\cap B(H)$, and
continuity gives the inclusion in $A(H)$.  The source uses etaleness only in
its subsequent attempt to identify $A_\sigma(H)$ with $A(H)$ by splitting
$X$ into individual source fibers.  We replace that step.

We shall use only the elementary coefficient-norm estimate
\[
 \|\pi_{\xi,\eta}\|_{B(H)}\leq \|\xi\|\,\|\eta\|.
\]

\section{Proof intuition}

The tempting proof disintegrates $\mu$ over $H\backslash X$ and identifies
$\sigma$ with an amplification of $\lambda$.  That route introduces
second-countability and measurable-selection hypotheses which are absent from
the question.

Instead, average a compactly supported function on $X$ along each $H$-orbit.
For $f\in C_c(X)$ this produces a bounded intertwiner
$W_f:L^2(H)\to L^2(X)$.  Properness of the free isotropy action gives the
bound.  More importantly, the whole family $\{W_f:f\in C_c(X)\}$ separates
the vectors of $L^2(H)$: restrict to one closed orbit and extend an arbitrary
$C_c(H)$ test function to $C_c(X)$.  Thus the adjoint ranges $\ran W_f^*$
span $L^2(H)$ densely.  A regular coefficient whose two vectors lie in finite
sums of these ranges is, by the intertwining equations, a finite sum of
$\sigma$-coefficients. Approximation in coefficient norm yields
$A(H)\subseteq A_\sigma(H)$, and \eqref{eq:source-reduction} finishes the
proof.

\section{The theorem}

\begin{theorem}[Full restriction theorem]\label{thm:main}
Let $\G$ be a locally compact Hausdorff groupoid admitting a left Haar system.
For every $u\in\G^{(0)}$, the restriction map
\[
           R_u:A(\G)\longrightarrow A(\G_u^u)
\]
is a contractive surjective algebra homomorphism.
\end{theorem}

The contractive homomorphism assertion is Proposition 3.6 of the source.  We
prove surjectivity through three lemmas.

\begin{lemma}[The isotropy action]\label{lem:proper}
The left action of $H$ on $X$ is free and proper.
\end{lemma}

\begin{proof}
If $hx=x$, multiplication on the right by $x^{-1}$ gives $h=u$, so the action
is free. Consider the shear map
\[
 \Theta:H\times X\longrightarrow X\times X,\qquad \Theta(h,x)=(hx,x).
\]
Its image is the equivalence relation
\[
 \mathcal R=\{(y,x)\in X\times X:s(y)=s(x)\}.
\]
The inverse on $\mathcal R$ is $(y,x)\mapsto(yx^{-1},x)$, so $\Theta$ is a
homeomorphism onto $\mathcal R$.  Since $\G^{(0)}$ is Hausdorff and $s$ is
continuous, $\mathcal R$ is closed in $X\times X$. Therefore $\Theta$ is
proper, which is precisely properness of the action.
\end{proof}

\begin{lemma}[A separating family of intertwiners]\label{lem:W}
For each $f\in C_c(X)$, the formula initially defined for $\xi\in C_c(H)$ by
\begin{equation}\label{eq:W}
       (W_f\xi)(x)=\int_H \xi(h)f(h^{-1}x)\,dh
\end{equation}
extends to a bounded operator $W_f:L^2(H)\to L^2(X,\mu)$ satisfying
\[
       W_f\lambda(k)=\sigma(k)W_f\qquad(k\in H).
\]
Moreover,
\[
       \bigcap_{f\in C_c(X)}\ker W_f=\{0\}.
\]
\end{lemma}

\begin{proof}
Set
\[
     a_f(x)=\int_H |f(h^{-1}x)|\,dh.
\]
Properness implies that $a_f$ is continuous: near a fixed $x$, the $h$ for
which $h^{-1}x\in\supp f$ lie in one compact subset of $H$, so ordinary
compact-support dominated convergence applies.  The function is
$H$-invariant and, on the quotient $H\backslash X$, its support is contained
in the compact set which is the image of $\supp f$. Consequently
$M_f:=\sup_X a_f<\infty$.

By Cauchy--Schwarz and Tonelli,
\begin{align*}
 |W_f\xi(x)|^2
 &\le a_f(x)\int_H |\xi(h)|^2|f(h^{-1}x)|\,dh,\\
 \|W_f\xi\|_2^2
 &\le M_f\int_H |\xi(h)|^2
       \left(\int_X|f(h^{-1}x)|\,d\mu(x)\right)dh\\
 &=M_f\|f\|_{L^1(X,\mu)}\|\xi\|_2^2.
\end{align*}
The last equality is exactly the left invariance of the Haar-system measure
$\mu=\lambda^u$. Thus $W_f$ extends boundedly. A change of variables
$h=k\ell$ gives
\[
 W_f\lambda(k)\xi(x)
 =\int_H\xi(\ell)f(\ell^{-1}k^{-1}x)\,d\ell
 =\sigma(k)W_f\xi(x).
\]

We next record a pointwise fact needed for separation. For every
$\xi\in L^2(H)$, $W_f\xi$ has a continuous representative. Indeed, on any
compact subset of $X$, properness confines the supports of
$h\mapsto f(h^{-1}x)$ to a common compact subset of $H$, and the map
$x\mapsto[h\mapsto f(h^{-1}x)]$ is continuous into $L^2(H)$. Approximation of
$\xi$ by $C_c(H)$ therefore gives locally uniform convergence of the
corresponding functions in \eqref{eq:W}.

Suppose now that $W_f\xi=0$ in $L^2(X,\mu)$ for every $f\in C_c(X)$.  A Haar
system has full support, so continuity forces each $W_f\xi$ to vanish
pointwise. Fix $x_0\in X$. Its orbit $O=Hx_0$ is closed, and the map
\[
       H\longrightarrow O,\qquad h\longmapsto h^{-1}x_0
\]
is a homeomorphism. Given $\phi\in C_c(H)$, define $g\in C_c(O)$ by
$g(h^{-1}x_0)=\phi(h)$. The standard compact-support extension lemma for
closed subspaces of locally compact Hausdorff spaces supplies
$f\in C_c(X)$ with $f|_O=g$.  For clarity, this lemma needs no normality
assumption on $X$: extend $g$ by $0$ at infinity from $O\cup\{\infty\}$ in
the compact Hausdorff one-point compactification of $X$, apply Tietze there,
and multiply by a compactly supported cutoff equal to $1$ on $\supp g$.
Then
\[
       0=(W_f\xi)(x_0)=\int_H\xi(h)\phi(h)\,dh.
\]
As $\phi$ was arbitrary and $C_c(H)$ is dense in $L^2(H)$ for the Radon Haar
measure, $\xi=0$. This proves the common-kernel assertion.
\end{proof}

\begin{lemma}[Separating intertwiners and coefficients]\label{lem:coefficients}
Let $\pi$ and $\rho$ be unitary representations of a locally compact group.
If a family of bounded intertwiners $T_i:\Hil_\pi\to\Hil_\rho$ has
$\bigcap_i\ker T_i=\{0\}$, then
\[
                         A_\pi(H)\subseteq A_\rho(H).
\]
\end{lemma}

\begin{proof}
The orthogonal complement of
$\overline{\operatorname{span}}_i\ran T_i^*$ is $\bigcap_i\ker T_i$, so the
former space is all of $\Hil_\pi$. It is therefore enough to consider
coefficient vectors
\[
       \xi'=\sum_{i=1}^mT_i^*\alpha_i,
       \qquad \eta'=\sum_{j=1}^nT_j^*\beta_j.
\]
For every group element $k$, the adjoint intertwining relation gives
\begin{align*}
 \langle\pi(k)\xi',\eta'\rangle
 &=\sum_{i,j}\langle\pi(k)T_i^*\alpha_i,T_j^*\beta_j\rangle\\
 &=\sum_{i,j}\langle\rho(k)T_jT_i^*\alpha_i,\beta_j\rangle.
\end{align*}
This is a finite sum of coefficients of $\rho$. Approximating arbitrary
$\xi,\eta\in\Hil_\pi$ by such $\xi',\eta'$ and using
\[
 \|\pi_{\xi,\eta}-\pi_{\xi',\eta'}\|_{B(H)}
 \leq \|\xi-\xi'\|\,\|\eta\|+\|\xi'\|\,\|\eta-\eta'\|
\]
proves that every coefficient of $\pi$ belongs to $A_\rho(H)$.
\end{proof}

\begin{proof}[Proof of Theorem~\ref{thm:main}]
Apply Lemmas \ref{lem:W} and \ref{lem:coefficients} with $\pi=\lambda$ and
$\rho=\sigma$. They give directly
\[
           A(H)=A_\lambda(H)\subseteq A_\sigma(H).
\]
Together with the source reduction \eqref{eq:source-reduction}, this gives
$A_\sigma(H)=A(H)$ and
\[
           R_u(A(\G))=A(H).
\]
This is surjectivity. Contractivity and multiplicativity were already
established in the source at the stated generality.
\end{proof}

\section{Scope, checks, and novelty}

\paragraph{Why the former extra hypotheses disappear.}
The argument never restricts $\lambda^u$ to $H$ and never assigns positive
measure to an individual orbit.  Thus it applies even when
$\lambda^u(H)=0$, as in nondiscrete pair groupoids.  It uses only left Haar
measure and left invariance, so nonunimodular isotropy groups cause no modular
factor.  Every integral is first taken over compact supports; no separability,
sigma-finiteness, Borel section, or disintegration theorem is used.

\paragraph{Verification.}
The proof was audited at its four nonformal points: (i) the shear-map proof of
properness; (ii) the Schur--Cauchy bound for $W_f$; (iii) continuity plus
compact-support extension from a closed orbit; and (iv) the adjoint-range
coefficient calculation in Lemma \ref{lem:coefficients}. No computational
check is relevant. The most valuable human check is a line-by-line comparison
of \eqref{eq:source-reduction} with Claim 1 in the source proof.

\paragraph{Bounded novelty search.}
On 27 August 2026, bounded searches used the exact arXiv id and title, the
authors' names, and the phrases ``Fourier algebra groupoid restriction
isotropy surjective'', ``$A(\G)\to A(\G_u^u)$'', and ``proper action Koopman
regular representation''.  The cached source, including its May 2026 update,
still contains Remark 3.10. The searches found the source and general
proper-action/Koopman background, but no later work stating this
countability-free restriction theorem. This is evidence, not a guarantee, of
novelty.

\paragraph{Human-review recommendation.}
Treat this as a likely-valid candidate full solution pending expert review.
In particular, verify the compact-support control in Lemma \ref{lem:W}, the
closed-orbit extension argument, and the source reduction. If those checks
pass, the result answers Remark 3.10 exactly.

\begin{thebibliography}{9}
\bibitem{DG}
J. DeGaetani and M. Ghandehari,
\emph{On the restriction maps of the Fourier and Fourier-Stieltjes algebras
over locally compact groupoids}, arXiv:2411.14624 (2024; updated 2026).

\bibitem{Arsac}
G. Arsac, \emph{Sur l'espace de Banach engendr\'{e} par les coefficients
d'une repr\'{e}sentation unitaire}, Publ. D\'{e}p. Math. (Lyon) 13 (1976),
no. 2, 1--101.

\bibitem{KL}
E. Kaniuth and A. T. Lau, \emph{Fourier and Fourier--Stieltjes Algebras on
Locally Compact Groups}, Mathematical Surveys and Monographs 231, American
Mathematical Society, 2018; see in particular Section 2.8.
\end{thebibliography}

\end{document}
